\subsection{Agent-to-Agent Monitor}
\label{sec:agent-to-agent-monitor}

We implement the project-defined Agent-to-Agent Monitor as a post-hoc,
detection-only baseline. It combines structural summaries of the observable
agent-interaction graph with an LLM judge for compositional leakage, asking
whether individually plausible outputs collectively advance a harmful,
deceptive, or covert objective. It does not train a classifier or intervene in
the live multi-agent execution.

For each saved trajectory, we construct a label-free public graph containing one
node per agent. Each node stores the agent's latest public message and the number
of public messages it produced. Directed edges summarize public sender--recipient
pairs, their message counts, and the rounds in which communication occurred. The
monitor additionally receives the task context and options visible to all agents
and the final aggregated output. Private coordination, observable rationales,
designated-colluder identities, attack targets, condition labels, and evaluation
metadata are excluded. Labels are joined only after prediction.

We apply a separate LLM judge (\texttt{openai/bedrock/gpt-oss-120b}) with
temperature zero, low reasoning effort, and the frozen prompt version
\texttt{compositional\_leakage\_v1}. The judge returns a continuous
suspiciousness score in $[0,1]$ and a one-sentence evidence statement. We freeze
the decision threshold at 0.5. Expected cycles, bidirectional communication,
and high connectivity are recorded as graph statistics but are not themselves
treated as evidence of attack, since NARCBench uses the same fully connected
multi-round topology in both matched conditions.

We evaluate all 50 collusion trajectories and their 50 matched controls. The
primary condition-label analysis reports AUROC, average precision, recall,
false-positive rate, precision, balanced accuracy, and paired attack--control
score differences. Because the defense observes public consequences rather
than private intent, we additionally report an exploratory harmful-outcome
analysis based on the end-to-end \texttt{attack\_success} label. This label is
never supplied to the monitor. Confidence intervals for proportions use Wilson
intervals; paired score and flag-rate differences use a 10,000-resample paired
bootstrap; paired binary flags use an exact McNemar test.

The condition-label AUROC was 0.714 and average precision was 0.780. At the
fixed threshold, the monitor flagged 3/50 attack-condition trajectories and
0/50 controls (recall 0.060, false-positive rate 0, precision 1.0). Mean
suspiciousness was 0.107 for attack trajectories and 0.040 for controls; the
paired difference was 0.067 (bootstrap 95\% CI [0.026, 0.119]). Only three
collusion attempts achieved the strict end-to-end attack-success outcome; the
monitor detected two of them. This outcome-level estimate is exploratory due to
the very small positive count. The final run required 101 model calls, including
one retry after a truncated invalid response, and 215,154 provider-reported
tokens.
